\documentclass[12pt]{article}
\usepackage{amsmath, amssymb, amsthm}
\usepackage{geometry}
\geometry{a4paper, margin=1in}

% Theorem and Lemma Environments
\newtheorem{theorem}{Theorem}[section]
\newtheorem{lemma}[theorem]{Lemma}
\newtheorem{proposition}[theorem]{Proposition}
\newtheorem{corollary}[theorem]{Corollary}
\theoremstyle{definition}
\newtheorem{definition}[theorem]{Definition}
\newtheorem{remark}[theorem]{Remark}

\begin{document}

\title{Mild Initial Conditions for Global Regularity}
\author{}
\date{}
\maketitle

\section{Mild Initial Conditions: $u_0 \in H^1$}

The global regularity of the Navier--Stokes equations begins with the assumption of mild initial conditions, specifically that the initial velocity field $u_0$ belongs to the Sobolev space $H^1(\mathbb{R}^3)$:
\begin{equation}
\|u_0\|_{H^1} = \|u_0\|_{L^2} + \|\nabla u_0\|_{L^2} < \infty.
\end{equation}

This assumption ensures finite energy and enstrophy at the outset, providing a foundation for energy estimates, spectral decay, and compactness arguments.

\subsection{Properties of $H^1$ Initial Conditions}

\paragraph{Finite Energy and Enstrophy:}
The Sobolev norm $\|u_0\|_{H^1}$ captures:
\begin{itemize}
    \item **Kinetic Energy:** $\|u_0\|_{L^2}^2 = \int_{\mathbb{R}^3} |u_0(x)|^2 \, dx$.
    \item **Enstrophy:** $\|\nabla u_0\|_{L^2}^2 = \int_{\mathbb{R}^3} |\nabla u_0(x)|^2 \, dx$.
\end{itemize}
Finite enstrophy is critical for controlling velocity gradients, ensuring smoothness of solutions at small scales and aligning with spectral decay arguments.

\paragraph{Compatibility with Weak Solutions:}
Initial conditions in $H^1$ are compatible with the weak formulation of the Navier--Stokes equations. For $u_0 \in H^1$, there exists a weak solution $u \in L^2(0, T; H^1)$ satisfying:
\begin{equation}
\int_0^T \int_{\mathbb{R}^3} \left( \partial_t u \cdot v + (u \cdot \nabla) u \cdot v + \nabla p \cdot v - \nu \Delta u \cdot v \right) \, dx \, dt = 0,
\end{equation}
for all test functions $v \in H^1$.

\paragraph{Extension to Strong Solutions:}
While $H^1$ initial conditions suffice for weak solutions, regularity results require additional arguments (e.g., energy decay, compactness) to extend these solutions to strong solutions. Finite enstrophy plays a critical role in controlling high-frequency energy, preventing singularities.

\subsection{Connection to Energy Estimates and Spectral Decay}
The initial energy and enstrophy, represented by $\|u_0\|_{H^1}$, propagate through the evolution of the monotone functional $Q(t)$:
\begin{equation}
Q(t) = \|u(t)\|_{L^2}^2 + \|\nabla u(t)\|_{L^2}^2.
\end{equation}
For $u_0 \in H^1$, $Q(0)$ is finite, ensuring that energy estimates hold globally.

\begin{lemma}[Propagation of Energy and Enstrophy]
If $u_0 \in H^1$, then for all $t \geq 0$:
\begin{equation}
Q(t) \leq Q(0) e^{-\nu t} + \int_0^t \|f(s)\|_{L^2}^2 \, ds,
\end{equation}
where $f$ represents external forcing (if present).
\end{lemma}

\begin{proof}
The result follows from integrating the energy inequality:
\begin{equation}
\frac{dQ}{dt} + 2\nu \|\nabla u\|_{L^2}^2 \leq \|f\|_{L^2} \|u\|_{L^2}.
\end{equation}
By Gronwall’s inequality, $Q(t)$ remains bounded for all $t \geq 0$.
\end{proof}

The boundedness of $Q(t)$ directly supports the spectral decay assumption:
\begin{equation}
E(k, t) \leq C k^{-\beta}, \quad \beta > 3,
\end{equation}
ensuring smooth energy transfer and avoiding unbounded growth at high wave numbers.

\subsection{Role in Compactness and Regularity}
The $H^1$ condition provides the foundation for compactness arguments by ensuring:
\begin{itemize}
    \item **Spectral Decay:** High-frequency energy contributions decay sufficiently fast, supporting the Scale-Barrier Principle.
    \item **Strong Convergence:** Bounded initial energy propagates to strong convergence in $L^2$ via Non-Sobolev Compactness.
    \item **Gradient Control:** Finite enstrophy enables curvature-based damping of velocity gradients, ensuring smoothness.
\end{itemize}

\subsection{Relevance to Global Regularity}
Mild initial conditions are essential for global regularity because:
\begin{itemize}
    \item They ensure well-posedness of the Navier--Stokes equations in the weak sense.
    \item Finite enstrophy and energy propagate forward in time, preventing blowup.
    \item They align with the assumptions of compactness, spectral decay, and curvature-based regularity.
\end{itemize}

\section*{Conclusion}
Mild initial conditions in $H^1$ provide the essential starting point for establishing global regularity. By ensuring finite energy and enstrophy, they facilitate energy estimates, spectral decay, compactness, and gradient control, forming the basis for the subsequent steps in the unified proof.

\section*{References}
\begin{itemize}
    \item O. Ladyzhenskaya, \emph{The Mathematical Theory of Viscous Incompressible Flow}, Gordon and Breach, 1969.
    \item R. Temam, \emph{Navier--Stokes Equations: Theory and Numerical Analysis}, North-Holland, 1977.
    \item H. Triebel, \emph{Theory of Function Spaces}, Birkh\"auser, 1983.
    \item P. Constantin, "Geometric Methods in Fluid Dynamics," Comm. Math. Phys., 1990.
\end{itemize}

\end{document}
